\subsection{极大秩子群}

我们把连通李群 G 的李代数的秩称为 G 的秩，记作 rk G。由定理 2 a)，连通紧李群的秩是其诸极大环面的公共维数。

设 G 是连通紧李群，H 是 G 的闭子群。若 H 连通，则 \(rkH \leq rkG\) （因为 H 的极大环面是 G 中的环面）。由定理 2 c)，说 H 连通且为极大秩（即为秩 rkG），就等于说 H 是 G 的某些极大环面之并。我们立即由命题 1 推出：

命题 2. 设 \(f: \mathrm{G} \to \mathrm{G}'\) 是核为中心的连通紧李群满态射。映射 \(\mathrm{H} \mapsto f(\mathrm{H})\) 与 \(\mathrm{H}' \mapsto f^{-1}(\mathrm{H}')\) 是 \(\mathrm{G}\) 的极大秩连通闭子群集与 \(\mathrm{G}'\) 的相应子群集之间的互逆双射。

命题 3. 设 G 是连通紧李群，H 是极大秩连通闭子群。

\begin{enumerate}
\def\labelenumi{\alph{enumi})}
\item
  紧流形 G/H 是单连通的。
\item
  由 H 到 G 中的典则单射诱导的同态 \(\pi_1(\mathrm{H})\to \pi_1(\mathrm{G})\) 是满射。
\end{enumerate}

由于 H 连通，我们有正合列（《一般拓扑学》，第 XI 章，编写中）

\[
\pi_ {1} (\mathrm{H}) \rightarrow \pi_ {1} (\mathrm{G}) \rightarrow \pi_ {1} (\mathrm{G/H}, \bar {e}) \rightarrow 0
\]

其中 \(\bar{e}\) 是 G 的单位元在 G/H 中的像。由于 G/H 连通，这立即给出断言 a) 与 b) 的等价性。此外，若 \(f : G' \to G\) 是核为中心的连通紧李群满态射，则对 G 证明本命题（取 a) 的形式）与对 \(G'\) 证明它是一回事（命题 2）。于是，我们可以先把 G 换成 \(\text{Ad}(G)\) ，再假定 G 半单，然后通过把 G 换成一个泛覆盖（ \(\S1\) ，第 4 节，命题 4 的推论 2），假定 G 单连通。但此时断言 b) 是平凡的。

命题 4. 设 G 是紧李群，H 是 G 的极大秩连通闭子群，N 是 H 在 G 中的正规化子。则 H 在 N 中指数有限，且是 N 的单位元连通分支。

事实上，H 的李代数包含 L(G) 的一个嘉当子代数。于是，由第 VII 章，§2，第 1 节，命题 4 的推论 4，H 是 N 的单位元连通分支。由于 N 紧，H 在 N 中指数有限。

注。1) 满足 rkH = rkG 的 G 的每个积分子群 H 都是闭的：事实上，前面的证明表明 H 是其正规化子的单位元连通分支，而该正规化子是 G 的闭子群。

\begin{enumerate}
\def\labelenumi{\arabic{enumi})}
\setcounter{enumi}{1}
\tightlist
\item
  采用命题 4 的记号，每个闭子群 \(H'\) （G 中包含 H 且使 \((H':H)\) 有限者）都正规化 H，故含于 N；同样，\(H'\) 的正规化子也含于 N。特别地，N 是自身的正规化子。
\end{enumerate}

